\section{Source audit and proposed corrections}
\label{sec:audit}

The inspected sources support the preceding developments, but the
research-status map needs an additional cross-reference.  The correction
below concerns attribution and an already answered construction problem.
It does not invalidate the same-direction spectral representations.
No counterexample to the inspected source reconstruction or quadratic
factorization theorems was found.

\subsection{The connected double family is an established interpolant}

For $\Re q>0$ write
\[
 Z_q(s_1,\ldots,s_d)=
  \sum_{n_1>\cdots>n_d\ge0}\prod_{j=1}^d(n_j+q)^{-s_j},
 \qquad Z_q(\varnothing)=1,
\]
initially in an absolute-convergence domain.  The same-direction report
defines
\begin{equation}\label{audit:connected}
 \mathcal C_{A,B}(q)=\zeta(A,B)-Z_q(A,B)
                      +\zeta(A,q)\{\zeta(B,q)-\zeta(B)\}.
\end{equation}
The earlier independent-orders report defines the relative interpolant
by ordered convolution \cite{IndependentOrders}:
\begin{equation}\label{audit:relative}
 Z_b(w)=\sum_{w=uv}Z_a(u)\mathcal H_{a,b}(v),
 \qquad \mathcal H_{a,b}(\varnothing)=1.
\end{equation}
The sum is over prefix--suffix decompositions of the word.  At depth two,
\begin{align*}
 \mathcal H_{q,1}(A,B)
 &=\zeta(A,B)-Z_q(A,B)
                  -\zeta(A,q)\{\zeta(B)-\zeta(B,q)\}.
\end{align*}
Thus there is an exact identity of the full functions,
\begin{equation}\label{audit:identification}
 \boxed{\mathcal C_{A,B}(q)=\mathcal H_{q,1}(A,B)
                              =\Psi(A,B;q-1).}
\end{equation}
The final member is the Ihara--Nakamura--Yamamoto interpolant
\cite[Proposition 2.1 and Corollary 4.2]{INY}.  Their theorem already
gives entire dependence on all independent indices at arbitrary depth.
The relative report explicitly cites that result and develops its
two-endpoint form.  This identification is algebraic on an open
convergence domain, so it does not rely on uniqueness from values at
integer $q$.

For clarity, the established general finite expression is
\begin{equation}\label{audit:star}
 \mathcal H_{q,1}(s_1,\ldots,s_d)=
  \sum_{j=0}^d(-1)^j
   Z_q^\star(s_j,\ldots,s_1)\,
                         \zeta(s_{j+1},\ldots,s_d).
\end{equation}
The empty prefix and suffix have value one, and $Z_q^\star$ uses weak
inequalities.  This is the convolution inverse of the strict ordered
series.  It preserves the published entire normalization, rather than
choosing separate finite parts for its meromorphic constituents.

The corresponding divided-difference kernel is also already available
\cite{INY,HarmonicGaps}.  If $X_0=0$, $X_j=v_1+\cdots+v_j$, then an
equivalent convenient form is
\begin{equation}\label{audit:kernel}
 K_q^{(d)}(v_1,\ldots,v_d)
   =(-1)^d[e^{X_0},e^{X_1},\ldots,e^{X_d}]
                         (y\mapsto y^{d-q}).
\end{equation}
Here the brackets are the divided difference in $y$, extended
confluently when nodes coincide.  Its normalized Mellin transform is
$\mathcal H_{q,1}$, initially when every $\Re s_j>0$.
At $d=2$ it is precisely the kernel of the same-direction report.
To compare with the earlier formula using nodes $e^{-X_j}$, invert
each node in its finite barycentric sum: the product of all nodes
cancels the exponential prefactor, leaving the sign $(-1)^d$ and
the power $d-q$ in \eqref{audit:kernel}.

The shift law and the depth-two stuffle law in this normalization are
\begin{align}
 \mathcal H_{q+1,1}(s_1,\ldots,s_d)
 &=\mathcal H_{q,1}(s_1,\ldots,s_d)
                   +q^{-s_1}\mathcal H_{q,1}(s_2,\ldots,s_d),
                   \label{audit:shift}\\
 h_q(s)h_q(t)
 &=\mathcal H_{q,1}(s,t)+\mathcal H_{q,1}(t,s)+h_q(s+t),
 \quad h_q(s)=\zeta(s)-\zeta(s,q).
 \label{audit:stuffle}
\end{align}
They are inherited identities for entire functions.  Sign conventions
can be checked immediately: the same-direction report uses the
depth-one difference $\zeta(s,q)-\zeta(s)=-h_q(s)$, which accounts
for the apparent change of sign in its displayed laws.

\paragraph{Proposed source edit.}
The attribution paragraph of \cite{SameDirection} should identify its
entire connected double family as the depth-two specialization of
\cite{INY,IndependentOrders}.  Its second further-research question
should no longer ask for existence, entireness, shift/stuffle laws, or
the basic higher divided-difference kernel.  Those constructions are
already available.  A useful replacement asks for explicit minimal
coordinate reductions on specified mixed nonpositive-index loci,
or for comparisons between genuinely different normalizations.
The Barnes and beta--Newton same-direction representations remain
separate results and are not affected by this correction.

\subsection{Updates justified by the new proofs}

The order-eight statement in \cite{Reconstruction} correctly retained
the congruence qualification $b\not\equiv1\pmod{16}$.  It was not a
false theorem.  After the present proof and certificate are accepted,
the collective statement may be strengthened to every $b\ge5$ and its
Q1 closed.  The original congruence proof remains useful as the finite
range reduction.  Its previously rejected simple valuation extrapolation
should remain rejected; the new proof does not restore that formula.

The same report's Q5 is now answered for the fourth power and its
regular even moments.  The statement should refer to the six displayed
formal coordinates and the all-index Bernoulli transform.  It should
not assert independence of the evaluated constants or promote fifth
and higher powers to the same status.  The starting value $V_0=Q_7$
must retain Hoffman's attribution.

For \cite{Quadratic}, Q7 is answered by the pairwise logarithmic
correction and the confluent resolvent formulas, while Q2 is answered
by Corollary~\ref{cor:linear-jets}.  The latter uses the source's
linear factors and spectral scale two.  A product of two quadratic
factors has scale four; confusing them changes every higher correction.
The explicit primitive in Theorem~\ref{thm:primitive} is progress on
the mixed primitive direction, not a classification of all shift--mass
integrals.

\subsection{Corrections that must accompany higher-order reuse}

Three points are essential when extending the existing formulas.
First, the identity $\mathcal H_a(-m)=0$ does not imply
$\mathcal H_a'(-m)=0$.  At bilateral third order the omitted term
would be exactly $-12(B-C)^2\log(2\sin\pi a)$ for $P(v)=v^2$.
Second, multiplication by a Hurwitz pole uses one more coefficient
of the vanishing factor than multiplication by its regular part.
The term $A_{k,m+1}/(2r(m+1))$ in
\eqref{one:resonantC} is therefore required.  Its first additional
second-order example is $-\delta^2(B+C)/4$ in
\eqref{one:x5second}.  Third, an antiderivative of a summand is not
automatically summable: the constant in $\Phi_{B,C}$ is fixed so
that it decays to zero.  Without it, the separate series of primitives
would diverge.

These are requirements of the new extensions, not errors falsely
attributed to source theorems that only treated first derivatives.
The article and integration notes retain this distinction.

\subsection{Unchanged arithmetic conjectures}

The canonical Gaussian $S_6$ identity and the revised $S_8$ candidate
remain conjectural in this continuation \cite{Manuscript}.  The rejected
older $S_8$ vector remains different from the revised candidate.  No
failed relation search, small residual, or finite formal rank calculation
is converted here into an arithmetic proof or a nonreduction assertion.
Likewise, no new claim is made about the transcendence or numerical
independence of the six quartic coordinates.  These boundaries should
remain visible when the results are integrated into the collective text.
